% GENERATED FILE. Edit canonical lessons, then run npm run generate:books.
% canonical-source-hash: fnv1a64:13342dfa3b269e6a
% canonical-lessons: ZH-C156-yuan, ZH-W156-yuan-money, ZH-C156-jiaqian, ZH-C156-mianfei, ZH-C156-guanggao, ZH-C156-reading, ZH-C156-practice

\chapter{Price List}
\label{ch:zh-price-list}
\hlchaptermodality{156}

\begin{chapteropening}
\textbf{By the end of this chapter:} \emph{I can find prices in a price list and an advert.}

Reading a shop's price list and an advert, prices written in \zh{元} and said in kuài, the price, free of charge and an advert, built on one new written character: \zh{元}.
\end{chapteropening}

\section[\texorpdfstring{\zh{元} (yuán)}{元 (yuán)}]{\zh{元} (yuán) — yuan, as written}

\label{lesson:ZH-C156-yuan}



\begin{quote}
A new chapter, and it opens on how much things cost.

\begin{itemize}
\raggedright
  \item \emph{Recall:} say \zh{火车十点半到} — \emph{huǒchē shí diǎn bàn dào} — \textbf{R1}, the chapter you have finished
  \item \emph{Recall:} say \emph{guānmén}, to close — \textbf{R2}
  \item \emph{Recall:} say \emph{kuài}, yuan, as said — \textbf{R3}, durable
  \item \emph{Recall:} say \zh{多少} — \emph{duōshao}, how much — \textbf{R3}
\end{itemize}
\end{quote}

\subsection*{You'll want to know: \zh{元}}
\textbf{\zh{元}} — \emph{yuán} — \textquotedblleft{}yuan\textquotedblright{}, the unit of money, as it is written.

\begin{quote}
\textbf{\zh{一本书十元}} \emph{yì běn shū shí yuán} — A book costs ten yuan.
\end{quote}

On a price list it is written \zh{元}; aloud, most people say \emph{kuài}: \emph{shí kuài}. The next lesson puts the character in your hand.

\begin{cousinweb}
\textbf{\zh{元}} (\emph{yuán}) is inherited from Middle Chinese \emph{ngjwon} (Guangyun, in Baxter's transcription), glossed 'great; the beginning; the head'. The unit of money is a modern sense, assumed here, not quoted.
\end{cousinweb}

\begin{grammarlens}[title={What you've built}]
The number, then \zh{元}: \zh{十元}. A thing's price follows the thing: \zh{一本书}, \zh{十元}.

When you say it, swap in \emph{kuài}: \emph{yì běn shū shí kuài}. \zh{一} falls before the dip of běn: \emph{yì běn}.
\end{grammarlens}

\subsection*{Your turn}
\begin{itemize}
\raggedright
  \item \emph{Say it:} \emph{yuán}
  \item \emph{Say it:} \emph{yuán}, once more
  \item \emph{Say it:} \zh{一本书十元} — \emph{yì běn shū shí yuán} — A book costs ten yuan.
  \item \emph{Say it:} \emph{yì běn shū shí kuài} — the same price, as people say it
  \item \emph{Recall:} say \emph{kuài}, then \zh{元}
\end{itemize}

\begin{tcolorbox}[breakable,colback=teal!4,colframe=teal!35!black,title={Before you move on}]
What does \zh{元} mean? (\textbf{Yuan, as written}.) How do people say it aloud? (\emph{\textbf{kuài}}.)
\end{tcolorbox}
\section[\texorpdfstring{\zh{元} (yuán)}{元 (yuán)}]{\zh{元} — \zh{二}, then \zh{儿} under it}

\label{lesson:ZH-W156-yuan-money}



\begin{quote}
The word first, then the character inside it.

\begin{itemize}
\raggedright
  \item \emph{Recall:} say \emph{yuán} — \textbf{R1}, one lesson back, before it decays
\end{itemize}

\begin{quote}
\zh{二} \zh{儿}
\end{quote}

\emph{Read it:} each one aloud, with nothing said first — \emph{èr}, \emph{ér}
\end{quote}

\subsection*{Script — \zh{二}, then \zh{儿} under it}
\hlblockfigure{\detokenize{figures/ZH-W156-yuan-money-filmstrip.pdf}}{How \zh{元} is written, stroke by stroke}

\begin{quote}
\zh{元}
\end{quote}

\textbf{Two characters you own, one above the other.}

\textbf{\zh{二}, two strokes}: the short upper bar, then the longer lower bar.

\textbf{\zh{儿}, two strokes}, under them: the left-falling stroke, then the vertical that bends right along the foot and hooks up.

2 + 2 = four.

\subsection*{Your turn}
\begin{itemize}
\raggedright
  \item \emph{Write it:} \zh{元}, all four, with the model in front of you
  \item \emph{Cover:} write it again from memory, and count the pieces, not the strokes
\end{itemize}

\begin{tcolorbox}[breakable,colback=teal!4,colframe=teal!35!black,title={Before you move on}]
How many strokes in \zh{元}? (\textbf{Four}.) Which character is on top? (\textbf{\zh{二}}.) And under it? (\textbf{\zh{儿}}.)

Source: \href{https://raw.githubusercontent.com/chanind/hanzi-writer-data/68d10a4b21150cae5e1ebbd223eed289cf32d90c/data/\%E5\%85\%83.json}{Hanzi Writer Data \zh{元}, pinned stroke paths and medians}.
\end{tcolorbox}
\section[jiàqián]{jiàqián — price}

\label{lesson:ZH-C156-jiaqian}



\begin{quote}
Before the new one: read the last lesson's character in a price, then two words for money.

\begin{itemize}
\raggedright
  \item \emph{Read it:} \zh{十元}, aloud — \emph{shí yuán}
  \item \emph{Recall:} say \emph{qián}, money — \textbf{R3}, durable
  \item \emph{Recall:} say \emph{guì}, expensive — \textbf{R3}
\end{itemize}
\end{quote}

\subsection*{You'll want to know: jiàqián}
\emph{jiàqián} — \textquotedblleft{}price\textquotedblright{}.

\begin{quote}
\emph{jiàqián tài guì le} — The price is too high.
\end{quote}

Its characters come later; for now it is a word for the ear.

\begin{cousinweb}
\emph{jiàqián} is built from jià 'price' and \emph{qián} 'money', the word you know. Both are inherited from Middle Chinese, \emph{kaeH} and \emph{dzjen} (Guangyun, in Baxter's transcription), glossed 'price; value' and 'coin'. The jià of \emph{fàngjià}, a holiday, sounds the same and is a different word.
\end{cousinweb}

\begin{grammarlens}[title={What you've built}]
\emph{jiàqián} names the price as a fact, the figure on the list; \emph{duōshao qián} asks it aloud.

\zh{太} … \zh{了} is still the complaint: \emph{jiàqián tài guì le}.
\end{grammarlens}

\subsection*{Your turn}
\begin{itemize}
\raggedright
  \item \emph{Say it:} \emph{jiàqián}
  \item \emph{Say it:} \emph{jiàqián}, once more
  \item \emph{Say it:} \emph{jiàqián tài guì le} — The price is too high.
  \item \emph{Recall:} say \emph{qián}, then \emph{jiàqián}
\end{itemize}

\begin{tcolorbox}[breakable,colback=teal!4,colframe=teal!35!black,title={Before you move on}]
What does \emph{jiàqián} mean? (\textquotedblleft{}Price\textquotedblright{}.) How do you ask a price aloud? (\emph{duōshao qián}.)
\end{tcolorbox}
\section[miǎnfèi]{miǎnfèi — free of charge}

\label{lesson:ZH-C156-mianfei}



\begin{quote}
Before the new one: last lesson's word, then the word for buying.

\begin{itemize}
\raggedright
  \item \emph{Recall:} say \emph{jiàqián}, price — \textbf{R1}, one lesson back
  \item \emph{Recall:} say \zh{买} — \emph{mǎi}, to buy — \textbf{R3}, durable
\end{itemize}
\end{quote}

\subsection*{You'll want to know: miǎnfèi}
\emph{miǎnfèi} — \textquotedblleft{}free (of charge)\textquotedblright{}.

\begin{quote}
\emph{shuǐ miǎnfèi} — The water's free.
\end{quote}

Its characters come later; for now it is a word for the ear.

\begin{cousinweb}
\emph{miǎnfèi} is built from miǎn 'to be spared' and fèi 'cost'. Both are inherited from Middle Chinese, \emph{mjenX} and \emph{phj+jH} (Guangyun, in Baxter's transcription), glossed 'to escape; to remove' and 'to use up; to spend'. Spared the spending: free.
\end{cousinweb}

\begin{grammarlens}[title={What you've built}]
\emph{miǎnfèi} follows the thing, as a price does: \emph{shuǐ miǎnfèi}, where a price list would say \zh{水} \zh{一元}.

Its first syllable dips: \emph{miǎn}.
\end{grammarlens}

\subsection*{Your turn}
\begin{itemize}
\raggedright
  \item \emph{Say it:} \emph{miǎnfèi}
  \item \emph{Say it:} \emph{miǎnfèi}, once more
  \item \emph{Say it:} \emph{shuǐ miǎnfèi} — The water's free.
  \item \emph{Recall:} say \zh{买}, then \emph{miǎnfèi}
\end{itemize}

\begin{tcolorbox}[breakable,colback=teal!4,colframe=teal!35!black,title={Before you move on}]
What does \emph{miǎnfèi} mean? (\textquotedblleft{}Free of charge\textquotedblright{}.) Where does it go? (\textbf{After the thing, as a price does}.)
\end{tcolorbox}
\section[guǎnggào]{guǎnggào — advert}

\label{lesson:ZH-C156-guanggao}



\begin{quote}
Before the new one: last lesson's word, then a thing worth advertising.

\begin{itemize}
\raggedright
  \item \emph{Recall:} say \emph{miǎnfèi}, free — \textbf{R1}, one lesson back
  \item \emph{Recall:} say \zh{手机} — \emph{shǒujī}, a phone — \textbf{R3}, durable
\end{itemize}
\end{quote}

\subsection*{You'll want to know: guǎnggào}
\emph{guǎnggào} — \textquotedblleft{}advert; advertisement\textquotedblright{}.

\begin{quote}
\emph{zhè shì shǒujī de guǎnggào} — This is an advert for a phone.
\end{quote}

Its characters come later; for now it is a word for the ear.

\begin{cousinweb}
\emph{guǎnggào} is built from guǎng 'wide' and gào 'to tell'. Both are inherited from Middle Chinese, \emph{kwangX} and \emph{kawH} (Guangyun, in Baxter's transcription), glossed 'broad' and 'to report': telling widely. 'An advert' is the pair's modern sense.
\end{cousinweb}

\begin{grammarlens}[title={What you've built}]
\zh{的} joins the thing to the advert: \emph{shǒujī de guǎnggào}, an advert for a phone.

An advert often has a price and a day on it, and little else.
\end{grammarlens}

\subsection*{Your turn}
\begin{itemize}
\raggedright
  \item \emph{Say it:} \emph{guǎnggào}
  \item \emph{Say it:} \emph{guǎnggào}, once more
  \item \emph{Say it:} \emph{zhè shì shǒujī de guǎnggào} — This is an advert for a phone.
  \item \emph{Recall:} say \emph{miǎnfèi}, then \emph{guǎnggào}
\end{itemize}

\begin{tcolorbox}[breakable,colback=teal!4,colframe=teal!35!black,title={Before you move on}]
What does \emph{guǎnggào} mean? (\textquotedblleft{}An advert\textquotedblright{}.) How do you say an advert for a phone? (\emph{shǒujī de guǎnggào}.)
\end{tcolorbox}
\section[(reading)]{(reading) — a price list and an advert}

\label{lesson:ZH-C156-reading}



\begin{quote}
A shop's price list, then an advert under it.

\begin{itemize}
\raggedright
  \item \emph{Recall:} say \emph{guǎnggào}, an advert — \textbf{R1}, one lesson back
  \item \emph{Recall:} say \zh{白菜} — \emph{báicài}, cabbage — \textbf{R3}, durable
\end{itemize}

\begin{quote}
\zh{本子}
\end{quote}

\emph{Read it:} it aloud, with nothing said first — \emph{běnzi}
\end{quote}

\begin{tcolorbox}[breakable,skin=enhanced,colback=orange!6,colframe=orange!45!black,boxrule=0.5pt,arc=1mm,left=6pt,right=6pt,top=4pt,bottom=4pt,fonttitle=\bfseries,title={Read this}]
\begin{quote}
\zh{书} \zh{十元一本} \zh{本子} \zh{两元一本} \zh{大白菜} \zh{三元} \zh{水} \zh{一元} \zh{手机} \zh{九百九十九元}
\end{quote}

\begin{quote}
\zh{星期六和星期天} \zh{两本书十五元}
\end{quote}

\emph{Read it:} the list through once without stopping; do not translate as you go

\emph{Read it:} the advert again, and work out what two books cost at the weekend
\end{tcolorbox}

\begin{tcolorbox}[breakable,skin=enhanced,colback=orange!6,colframe=orange!45!black,boxrule=0.5pt,arc=1mm,left=6pt,right=6pt,top=4pt,bottom=4pt,fonttitle=\bfseries,title={Read this}]
\begin{quote}
\emph{shū, shí yuán yì běn} \emph{běnzi, liǎng yuán yì běn} \emph{dà báicài, sān yuán} \emph{shuǐ, yì yuán} \emph{shǒujī, jiǔbǎi jiǔshíjiǔ yuán}
\end{quote}

\begin{quote}
\emph{xīngqīliù hé xīngqītiān} \emph{liǎng běn shū shíwǔ yuán}
\end{quote}

A price list puts the thing first and its price after, with \zh{一本} for each one where the thing is counted by the book. \zh{一} falls before the rise of yuán: \emph{yì yuán}, as before běn. The advert changes only one thing: the price at the weekend.
\end{tcolorbox}

\subsection*{Your turn}
\begin{itemize}
\raggedright
  \item What do two books cost on Saturday? (\textbf{Fifteen yuan: \zh{两本书十五元}}, said \emph{shíwǔ kuài}.)
  \item What do a book and a notebook cost together? (\textbf{Twelve yuan}, said \emph{shí'èr kuài}.)
  \item Which costs one yuan? (\textbf{The water: \zh{水} \zh{一元}}.)
\end{itemize}

\begin{itemize}
\raggedright
  \item \emph{Say it:} the price list aloud, once, saying \emph{kuài} for each \zh{元}
\end{itemize}

\begin{tcolorbox}[breakable,colback=teal!4,colframe=teal!35!black,title={Before you move on}]
What is the dearest thing on the list? (\textbf{\zh{手机}}, 999 yuan.) What does a notebook cost? (\textbf{\zh{两元}}, said \emph{liǎng kuài}.) Which days does the advert cover? (\textbf{\zh{星期六和星期天}}.)
\end{tcolorbox}
\section[Practice]{\zh{一本书十元}}

\label{lesson:ZH-C156-practice}



\begin{quote}
The chapter closes here, at the counter, asking about the list you have read.

\begin{itemize}
\raggedright
  \item \emph{Recall:} say \zh{水} \zh{一元} — \emph{shuǐ, yì yuán} — \textbf{R1}, the reading
  \item \emph{Recall:} say \emph{miǎnfèi}, free — \textbf{R2}
  \item \emph{Recall:} say \emph{jiàqián}, price — \textbf{R2}
\end{itemize}
\end{quote}

\subsection*{The exchange}
\begin{quote}
\textbf{A:} \emph{zhè běn shū duōshao qián} — How much is this book?
\end{quote}

\begin{quote}
\textbf{B:} \zh{一本书十元}
\end{quote}

\begin{quote}
\emph{yì běn shū shí yuán} — Ten yuan a book.
\end{quote}

\begin{quote}
\textbf{A:} \emph{jiàqián tài guì le} — That's too dear.
\end{quote}

\begin{quote}
\textbf{B:} \zh{星期六和星期天} \zh{两本书十五元}
\end{quote}

\begin{quote}
\emph{xīngqīliù hé xīngqītiān, liǎng běn shū shíwǔ yuán} — At the weekend, two books for fifteen.
\end{quote}

\begin{quote}
\textbf{A:} \emph{nǐ kàndào guǎnggào le ma} — Did you see the advert?
\end{quote}

\begin{quote}
\textbf{B:} \emph{kàndào le, shuǐ miǎnfèi} — Yes, and the water's free.
\end{quote}

\begin{quote}
\textbf{A:} \zh{本子呢}
\end{quote}

\begin{quote}
\emph{běnzi ne} — And notebooks?
\end{quote}

\begin{quote}
\textbf{B:} \zh{两元一本}
\end{quote}

\begin{quote}
\emph{liǎng yuán yì běn} — Two yuan each.
\end{quote}

The list says \zh{元} and the speakers read it as written; at a counter, most people would say \emph{kuài}. Both are right.

\begin{tcolorbox}[breakable,skin=enhanced,colback=orange!6,colframe=orange!45!black,boxrule=0.5pt,arc=1mm,left=6pt,right=6pt,top=4pt,bottom=4pt,fonttitle=\bfseries,title={Read this}]
\emph{Read it:} the exchange through once without stopping, then answer the question below
\end{tcolorbox}

\begin{grammarlens}[title={What you've built — reading prices, saying prices}]
Written in \zh{元}, said in \emph{kuài}:

\begin{itemize}
\raggedright
  \item \textbf{\zh{一本书十元}} \emph{yì běn shū shí yuán} — said \emph{yì běn shū shí kuài}.
  \item \textbf{\zh{两本书十五元}} \emph{liǎng běn shū shíwǔ yuán} — two books, fifteen yuan.
  \item \emph{jiàqián tài guì le} — the price is too high.
\end{itemize}

\textbf{Listen.} \emph{Yí ge shǒujī duōshao qián?}

\textbf{Speak.} Read three prices from a list you know, saying \emph{kuài} each time.

\emph{Write it:} \zh{十元} — \zh{元} is \zh{二} over \zh{儿}
\end{grammarlens}

\subsection*{Your turn}
\begin{itemize}
\raggedright
  \item \emph{Say it:} \emph{zhè běn shū duōshao qián} — How much is this book?
  \item \emph{Say it:} \zh{一本书十元} — \emph{yì běn shū shí yuán} — Ten yuan a book.
  \item \emph{Say it:} the whole exchange, both parts
  \item \emph{Write it:} \zh{一本书十元}
\end{itemize}

\begin{tcolorbox}[breakable,colback=teal!4,colframe=teal!35!black,title={Before you move on}]
Read ten yuan as written. (\textbf{\zh{十元}}.) Say it as people do. (\emph{shí kuài}.) Say the water is free. (\emph{shuǐ miǎnfèi}.)
\end{tcolorbox}
